
\documentclass[conference]{IEEEtran}
\IEEEoverridecommandlockouts

\usepackage{graphicx}
\usepackage{amsmath}
\usepackage{hyperref}

\begin{document}

\title{Humor Detection in Text}

\author{
\IEEEauthorblockN{Priyansh Kabra}
\IEEEauthorblockA{
MIT Bengaluru \\
B.Tech CSE-AI(B), 3rd Year \\
Email: ---
}
}

\maketitle

\begin{abstract}
This paper presents a study on humor detection in short texts using the Humor Detection Dataset from Kaggle. The objective is to identify whether a given sentence is humorous, leveraging semantic incongruity and contextual cues for accurate classification.
\end{abstract}

\begin{IEEEkeywords}
Humor Detection, NLP, Machine Learning, Semantic Incongruity, Text Classification
\end{IEEEkeywords}

\section{Introduction}
Humor is an integral part of human communication and presents a significant challenge for computational understanding. This project explores methods for detecting humor in text using natural language processing (NLP) and machine learning techniques.

\section{Dataset}
The project uses the ``Humor Detection Dataset'' available on Kaggle, consisting of over 200K short text samples labeled as humorous or non-humorous. These samples provide a diverse and extensive corpus for model training and evaluation.

\section{Methodology}
A supervised learning approach was used to detect humor in text. Various models were evaluated, and features related to semantics and contextual incongruity were extracted. The final model achieved a validation accuracy of 89.8\%.

\section{Results}
The model achieved a validation accuracy of 89.8\%, demonstrating its effectiveness in distinguishing humorous from non-humorous text.

\section{Real-World Applications}
Humor detection can be applied in several domains including chatbot design, social media monitoring, and sentiment analysis. It helps AI systems better understand user intent and create more natural and context-aware responses.

\section{Conclusion}
This work demonstrates the feasibility of automatic humor detection in text using NLP and machine learning. Future work can focus on cross-lingual humor detection and multimodal approaches incorporating images or speech.

\section*{References}
\begin{thebibliography}{00}
\bibitem{b1} Humor Detection Dataset, Kaggle, Available at: \url{https://www.kaggle.com/datasets/}, Accessed 2025.
\end{thebibliography}

\end{document}
